% !TeX root = ../教材题库_全部条目_转写工作稿.tex
% Source: XJ-B1-C06-EX-P239-Y0957
% 第6章 / 6.4 用样本估计总体 / 教材 p.239、240 / PDF 第246、247页
% OCR未校对：公式与文字均可能有误；未完成数学TeX转写；禁止直接发给学生
% 原题及相邻原解截图（核对用）：../题目原图/XJ-B1-C06-EX-P239-Y0957-01.webp, ../题目原图/XJ-B1-C06-EX-P239-Y0957-02.webp
% 本文件只在转写完成后作为题干使用；OCR 条目尚不是可靠题干。
\begin{bankitem}{6.4 例题 例题4}{第6章 / 6.4 用样本估计总体 / 教材 p.239、240 / PDF 第246、247页}{XJ-B1-C06-EX-P239-Y0957}{OCR未校对：公式与文字均可能有误；未完成数学TeX转写；禁止直接发给学生}
\begin{ocrdraft}
[教材p.239 / PDF第246页]
\par
0互假设通过抽样调查，获得了某市100户居民某年的月均用电量（单位：
\par
kW\ensuremath{\cdot}h)，按从小到大顺序排列为：
\par
13．2
\par
16．4
\par
21.5
\par
21.7
\par
22．1
\par
23．1
\par
24．2
\par
25．7
\par
26．4
\par
26．5
\par
27．3
\par
2\&5
\par
30．2
\par
32．1
\par
32．5
\par
33．2
\par
34．1
\par
30．4
\par
33．9
\par
34．3
\par
36．2
\par
37．2
\par
3\&7
\par
41.3
\par
41.6
\par
42．5
\par
44．0
\par
49．2
\par
50．0
\par
52．8
\par
52．8
\par
54．6
\par
55．6
\par
55．8
\par
56．4
\par
56．6
\par
57．8
\par
5\&1
\par
59．8
\par
60．5
\par
61.4
\par
61.6
\par
62．0
\par
62．4
\par
62．4
\par
62．4
\par
63．9
\par
64．3
\par
64．3
\par
65．9
\par
66．1
\par
6\&7
\par
6\&7
\par
69．7
\par
72．4
\par
72．8
\par
73．2
\par
74．9
\par
75．0
\par
75．5
\par
77．3
\par
77．8
\par
7\&8
\par
7\&9
\par
79．4
\par
79．9
\par
80．3
\par
82．2
\par
83．3
\par
83．6
\par
84．8
\par
85．1
\par
85．1
\par
87．8
\par
89．3
\par
90．5
\par
91.1
\par
91.5
\par
91.6
\par
93．3
\par
93．8
\par
94．8
\par
9\&3
\par
107．9
\par
109．0
\par
114．5
\par
93．9
\par
95．3
\par
96．3
\par
99．4
\par
115.4
\par
116．7
\par
11\&3
\par
119．2
\par
120．3
\par
120．5
\par
132．4
\par
139．6
\par
146．8
\par
176．2
\par
计算样本数据的中位数、众数，并据此估计该市每户居民月均用电量的中位数、
\par
众数．
\par
[教材p.240 / PDF第247页]
\par
解由于100为偶数，所以最中间存在两个数，第50位是65．9，第51位是
\par
65．9十66．1
\par
66．1，所以中位数为
\par
=66．0．由上述数据可知，众数为62．4．
\par
2
\par
因此，估计该市居民每户月均用电量的中位数是66．0，众数是62．4．
\par
众数、中位数和平均数均能反映数据的集中趋势，而它们作为一组数据的代表
\par
又具有不同的特点．我们应当根据问题的需要，选择合适的统计量来描述数据的集
\par
中趋势．
\end{ocrdraft}
\end{bankitem}
